(1) 三维自由玻色子激发态粒子数为 $N_{\rm ex}=(V/\lambda_T^3)g_{3/2}(z)$，其中 $z=e^{\beta\mu}\le1$。凝聚临界点化学势从下接近零，$g_{3/2}(1)=\zeta(3/2)$，因此最大激发密度 $n_{\rm ex,max}=\zeta(3/2)/\lambda_T^3$。
(2) 在临界点 $N/V=\zeta(3/2)(mk_BT_c/(2\pi\hbar^2))^{3/2}$，整理为 $T_c=[2\pi\hbar^2/(mk_B)] [N/(V\zeta(3/2))]^{2/3}$。量纲为开尔文；固定数密度下粒子越轻，临界温度越高。
(3) $T<T_c$ 时化学势取零，激发态容量为 $N_{\rm ex}=N(T/T_c)^{3/2}$，所以凝聚分数 $N_0/N=1-(T/T_c)^{3/2}$。忽略零能基态贡献，热内能 $U(T)=(3/2)Vk_BT\,\zeta(5/2)/\lambda_T^3$，来自能量密度积分的 $g_{5/2}(1)$。
(4) 因 $\lambda_T^{-3}\propto T^{3/2}$，$U\propto T^{5/2}$，对定体积求导 $C_V=(15/4)Vk_B\zeta(5/2)/\lambda_T^3$，故 $C_V\propto T^{3/2}$。这些式子要求三维自由理想气体热力学极限，有限体积使跃迁圆滑，相互作用会改变临界温度和临界附近涨落。
